\SourcePage{183}{186}
\section[Transitions produced by a perturbation]{Transitions produced by a perturbation
acting for a finite time}

Suppose that the perturbation $V(t)$ acts only during some finite time interval
(or, alternatively, that $V(t)$ decreases sufficiently rapidly as
$t\to\pm\infty$). Before the perturbation begins to act (or in the limit
$t\to-\infty$), let the system be in the $n$th stationary state of the
discrete spectrum. At any later instant its state is described by
\[
  \Psi=\sum_k a_{kn}\Psi_k^{(0)},
\]
where, to first order,
\begin{equation}
 \begin{aligned}
  a_{kn}=a_{kn}^{(1)}&=-\frac i\hbar\int_{-\infty}^t
    V_{kn}e^{i\omega_{kn}t}\,dt,\qquad k\ne n,\\
  a_{nn}=1+a_{nn}^{(1)}&=1-\frac i\hbar\int_{-\infty}^tV_{nn}\,dt.
 \end{aligned}
 \tag{41.1}
\end{equation}
\SourcePage{184}{187}
The integration limits in (40.5) have been selected so that every
$a_{kn}^{(1)}$ vanishes as $t\to-\infty$. Once the action of the perturbation
has ended (or as $t\to\infty$), the coefficients $a_{kn}$ acquire constant
values $a_{kn}(\infty)$, and the system is in the state whose wave function is
\[
  \Psi=\sum_k a_{kn}(\infty)\Psi_k^{(0)}.
\]
This function again satisfies the unperturbed wave equation, though it differs
from the original function $\Psi_n^{(0)}$. By the general rules, the squared
modulus of $a_{kn}(\infty)$ gives the probability that the system has energy
$E_k^{(0)}$, that is, that it is in the $k$th stationary state.

Thus a perturbation can carry the system from its original stationary state
into any other one. The probability of a transition from the initial ($i$th)
stationary state to the final ($f$th) stationary state is\footnote{For
uniformity, whenever transition probabilities are discussed below, we shall
denote the initial and final states by the indices $i$ and $f$, respectively.
We shall also write the indices of transition probabilities specifically in
the order $fi$, consistently with the convention for matrix-element indices.}
\begin{equation}
  w_{fi}=\frac1{\hbar^2}\left|
    \int_{-\infty}^{+\infty}V_{fi}e^{i\omega_{fi}t}\,dt
  \right|^2.
  \tag{41.2}
\end{equation}

Now consider a perturbation which, once it has arisen, continues to act
indefinitely (while remaining small throughout). In other words, $V(t)$ tends
to zero as $t\to-\infty$ and to a finite nonzero limit as $t\to\infty$.
Formula (41.2) cannot be used directly here, since its integral diverges. From
the physical standpoint, however, this divergence is immaterial and is easily
removed. Integration by parts gives
\[
 \begin{aligned}
  a_{fi}&=-\frac i\hbar\int_{-\infty}^tV_{fi}e^{i\omega_{fi}t}\,dt\\
  &=-\left.\frac{V_{fi}e^{i\omega_{fi}t}}{\hbar\omega_{fi}}
    \right|_{-\infty}^t
  +\int_{-\infty}^t\frac{\partial V_{fi}}{\partial t}
    \frac{e^{i\omega_{fi}t}}{\hbar\omega_{fi}}\,dt.
 \end{aligned}
\]
The first term vanishes at the lower limit; at the upper limit it formally
coincides with the expansion coefficients in (38.8). (The extra periodic
factor occurs simply because $a_{fi}$ are expansion coefficients of the full
wave function $\Psi$, whereas the $c_{fi}$ of \secref{38} are expansion coefficients
\SourcePage{185}{188}
of the time-independent function $\psi$.) It is consequently clear that the
limit of this term as $t\to\infty$ describes only the change of the original
wave function $\Psi_i^{(0)}$ caused by the ``constant part'' $V(+\infty)$ of
the perturbation, and therefore has nothing to do with transitions into other
states. The transition probability is instead given by the square of the
second term:
\begin{equation}
  w_{fi}=\frac1{\hbar^2\omega_{fi}^2}\left|
    \int_{-\infty}^{+\infty}\frac{\partial V_{fi}}{\partial t}
    e^{i\omega_{fi}t}\,dt
  \right|^2.
  \tag{41.3}
\end{equation}

The formulas obtained also apply when a transition proceeds from a discrete
state to a continuous-spectrum state. The only difference is that one then
speaks of the probability of transition from the specified ($i$th) state into
states for which the quantities $\nu_f$ (see the end of \secref{38}) lie between
$\nu_f$ and $\nu_f+d\nu_f$. Thus, for example, (41.2) must be written as
\begin{equation}
  dw_{if}=\frac1{\hbar^2}\left|
    \int_{-\infty}^{\infty}V_{fi}e^{i\omega_{fi}t}\,dt
  \right|^2d\nu_f.
  \tag{41.4}
\end{equation}

If $V(t)$ varies little over time intervals of order $1/\omega_{fi}$, the
integral in (41.2) or (41.3) is very small. In the limit of arbitrarily slow
variation of the applied perturbation, the probability of every transition
that changes the energy (that is, with nonzero $\omega_{fi}$) tends to zero.
Hence, when an applied perturbation changes sufficiently slowly
(\emph{adiabatically}), a system initially in a particular nondegenerate
stationary state continues to remain in that same state (see also \secref{53}).

In the opposite limit, when a perturbation is switched on very rapidly---that
is, \emph{suddenly}---the derivatives $\partial V_{fi}/\partial t$ become
infinite at the ``switching instant.'' In the integral of
$(\partial V_{fi}/\partial t)e^{i\omega_{fi}t}$, the comparatively slow factor
$e^{i\omega_{fi}t}$ may then be taken outside the integral at its value at
that instant. The remaining integration is immediate and gives
\begin{equation}
  w_{fi}=\frac{|V_{fi}|^2}{\hbar^2\omega_{fi}^2}.
  \tag{41.5}
\end{equation}
Transition probabilities for sudden perturbations can also be found when the
perturbation is not small.
\SourcePage{186}{189}
Let the system initially be in a state described by an eigenfunction
$\psi_i^{(0)}$ of the original Hamiltonian $\hat H_0$. If the Hamiltonian
changes suddenly (in a time short compared with the periods $1/\omega_{fi}$
of transitions from the given state $i$ to other states), the wave function
has no time to change and remains what it was before the perturbation. It is
no longer, however, an eigenfunction of the system's new Hamiltonian $\hat H$;
the state $\psi_i^{(0)}$ is therefore not stationary. By the general rules of
quantum mechanics, the probabilities $w_{fi}$ for transition into any of the
new stationary states are determined by expanding $\psi_i^{(0)}$ in the
eigenfunctions $\psi_f$ of $\hat H$:
\begin{equation}
  w_{fi}=\left|\int\psi_i^{(0)}\psi_f^*\,dq\right|^2.
  \tag{41.6}
\end{equation}

We now show how this general formula reduces to (41.5) when the change
$\hat V=\hat H-\hat H_0$ of the Hamiltonian is small. Multiply the equations
\[
  \hat H_0\psi_i^{(0)}=E_i^{(0)}\psi_i^{(0)},\qquad
  \hat H^*\psi_f^*=E_f\psi_f^*
\]
by $\psi_f^*$ and $\psi_i^{(0)}$, respectively, integrate over $dq$, and
subtract the two results term by term. Using also the self-adjointness of
$\hat H$, we obtain
\[
  (E_f-E_i^{(0)})\int\psi_f^*\psi_i^{(0)}\,dq
  =\int\psi_f^*\hat V\psi_i^{(0)}\,dq.
\]
If $\hat V$ is small, then to first order $E_f$ may be replaced by the nearby
unperturbed level $E_f^{(0)}$, and $\psi_f$ on the right-hand side by the
corresponding function $\psi_f^{(0)}$. It follows that
\[
  \int\psi_f^*\psi_i^{(0)}\,dq
  =\frac1{\hbar\omega_{fi}}
   \int\psi_f^{(0)*}\hat V\psi_i^{(0)}\,dq,
\]
so that (41.6) becomes (41.5).

\ProblemHeading{1.} A charged oscillator in its ground state is suddenly subjected to a uniform
electric field. Determine the probabilities for the oscillator to pass into
excited states under the action of this perturbation.

\SourcePage{187}{190}
\SolutionHeading Potential energy of an oscillator in a uniform field exerting
force $F$ upon it is
\[
  U(x)=\frac{m\omega^2}{2}x^2-Fx
  =\frac{m\omega^2}{2}(x-x_0)^2+\operatorname{const},
\]
where $x_0=F/m\omega^2$. Thus it again has a purely oscillator form, but with
its equilibrium position displaced. Consequently, the stationary-state wave
functions of the perturbed oscillator are $\psi_k(x-x_0)$, where $\psi_k(x)$
are the oscillator functions (23.12), while the initial wave function is
$\psi_0(x)$ from (23.13). Using these functions and expression (23.11) for the
Hermite polynomials, we find
\[
  \int_{-\infty}^{+\infty}\psi_0^{(0)}\psi_k\,dx
  =\frac{(-1)^k}{\sqrt{2^k\pi k!}}e^{-\xi_0^2/2}
   \int_{-\infty}^{\infty}e^{-\xi\xi_0}
   \frac{d^k}{d\xi^k}e^{-\xi^2+2\xi\xi_0}\,d\xi,
\]
where $\xi_0=x_0\sqrt{m\omega/\hbar}$. Integrating by parts $k$ times
reduces the integral here to
\[
  \xi_0^k\int_{-\infty}^{\infty}\exp(-\xi^2+\xi\xi_0)\,d\xi
  =\xi_0^k\sqrt\pi\exp\frac{\xi_0^2}{4}.
\]
Formula (41.6) therefore gives the required transition probability as
\[
  w_{k0}=\frac{\bar k^k}{k!}e^{-\bar k},\qquad
  \bar k=\frac{\xi_0^2}{2}=\frac{F^2}{2m\hbar\omega^3}.
\]
As a function of the integer $k$, this is a Poisson distribution with mean
$\bar k$.

The perturbative regime corresponds to small $F$ such that $\bar k\ll1$.
The excitation probabilities are then small and fall rapidly as $k$
increases; the largest is $w_{10}\approx\bar k$.

Conversely, when $F$ is large ($\bar k\gg1$), excitation occurs with
overwhelming probability: the chance that the oscillator remains in its
normal state is $w_{00}=e^{-\bar k}$.

\ProblemHeading{2.} The nucleus of an atom in its normal state receives a
sudden impulse and thereby acquires velocity $\mathbf v$. Assume that the
duration $\tau$ of the impulse is short both in comparison with the electronic
periods and with $a/v$, where $a$ is an atomic dimension. Determine the
probability that this ``shaking'' excites the atom (\emph{A.~B. Migdal}, 1939).

\SolutionHeading Pass to the reference frame $K'$ that moves with the nucleus
after the impact. Because $\tau\ll a/v$, the nucleus may be treated as having
hardly moved during the impact, and the electron coordinates in $K'$ and in
the original frame $K$ therefore coincide immediately after the perturbation.
The initial wave function in $K'$ is
\[
  \psi_0'=\psi_0\exp\left(-i\mathbf q\sum_a\mathbf r_a\right),
  \qquad \mathbf q=\frac{m\mathbf v}{\hbar},
\]
where $\psi_0$ is the normal-state wave function for a stationary nucleus and
the sum in the exponential extends over all $Z$ electrons of the atom.
\SourcePage{188}{191}
According to (41.6), the probability sought for transition into the $k$th
excited state is now
\[
  w_{k0}=\left|\left\langle k\left|
   \exp\left(-i\mathbf q\sum_a\mathbf r_a\right)\right|0\right\rangle\right|^2.
\]
In particular, for $qa\ll1$, expand the exponential under the integral. The
integral of $\psi_k^*\psi_0$ vanishes by orthogonality, and hence
\[
  w_{k0}=\left|\left\langle k\left|
    \mathbf q\sum_a\mathbf r_a\right|0\right\rangle\right|^2.
\]

\ProblemHeading{3.} Find the total probability of excitation and ionization
of a hydrogen atom subjected to a sudden ``shaking'' (see the preceding
problem).

\SolutionHeading The required probability can be evaluated as the difference
\[
  1-w_{00}=1-\left|\int\psi_0^2e^{-i\mathbf q\mathbf r}\,dV\right|^2,
\]
where $w_{00}$ is the probability that the atom remains in the ground state;
$\psi_0=(\pi a^3)^{-1/2}e^{-r/a}$ is the hydrogen ground-state wave function,
and $a$ is the Bohr radius. Evaluation of the integral gives
\[
  1-w_{00}=1-\frac1{(1+\tfrac14q^2a^2)^4}.
\]
In the limiting case $qa\ll1$, this probability tends to zero as
$1-w_{00}\approx q^2a^2$; for $qa\gg1$, it approaches unity as
$1-w_{00}\approx1-(2/qa)^8$.

\ProblemHeading{4.} Determine the probability that an electron is ejected
from the $K$ shell of an atom with large atomic number $Z$ during nuclear
$\beta$ decay. Assume the $\beta$-particle velocity to be large relative to
the velocity of a $K$ electron (\emph{A.~B. Migdal}, 1941;
\emph{E.~L. Feinberg}, 1939).

\SolutionHeading\footnote{Atomic units are used in Problems 4 and 5.} Under
the stated conditions, the time taken by the $\beta$ particle to cross the
$K$ shell is short compared with the electron's orbital period; the change in
nuclear charge may therefore be regarded as instantaneous. The perturbation
is the change $V=1/r$ in the nuclear field when its charge changes by an
amount small compared with $Z$ (namely, by 1). According to (41.5), the
probability for either of the two $K$-shell electrons, whose energy is
$E_0=-Z^2/2$,\footnote{Here and below the $K$-electron states are treated as
hydrogen-like (see \secref{74}).} to enter a continuous-spectrum state of energy
$E=k^2/2$ in the interval $dE=k\,dk$ is
\[
  dw=2\frac{4|V_{0k}|^2}{(k^2+Z^2)^2}\,dk.
\]
The integral defining $V_{0k}$ receives its essential contribution from
distances of order $1/Z$ near the nucleus. In this region a hydrogen-like
expression may also be used for the continuous-spectrum wave function. The
electron's final state must have $l=0$, the same angular momentum as the
initial state. From the function $R_{10}$ and the function $R_{k0}$ normalized
on the $k/2\pi$ scale, obtained in \secref{36},
\SourcePage{189}{192}
together with formula (f.3) of the Mathematical Appendices, we find\footnote{It
is convenient during the calculation to use Coulomb units, returning to
atomic units in the final result.}
\[
  \left(\frac1r\right)_{0k}
  =\frac{4\sqrt{2\pi k}}{1-e^{-2\pi Z/k}}
   \frac{(1+ik/Z)^{iZ/k}(1-ik/Z)^{-iZ/k}}{1+k^2/Z^2},
\]
and, since
\[
  |(1+i\alpha)^{i/\alpha}|^2
  =\exp\left(-2\frac{\arctan\alpha}{\alpha}\right),
\]
we finally obtain
\[
  dw=\frac{2^7}{Z^4(1+k^2/Z^2)^4}f\left(\frac kZ\right)k\,dk,
\]
where
\[
  f(\alpha)=\frac1{1-e^{-2\pi/\alpha}}
  \exp\left(-4\frac{\arctan\alpha}{\alpha}\right).
\]
The limiting values of $f(\alpha)$ are
\[
  f=e^{-4}\quad\text{for }\alpha\ll1,\qquad
  f=\alpha/2\pi\quad\text{for }\alpha\gg1.
\]
The total probability of ionizing the $K$ shell follows by integrating $dw$
over all energies of the emitted electron. Numerical evaluation yields
$w=0.65Z^{-2}$.

\ProblemHeading{5.} Determine the probability that an electron is ejected
from the $K$ shell of a large-$Z$ atom during nuclear $\alpha$ decay. The
$\alpha$ particle moves slowly compared with a $K$ electron, although its
time of emergence from the nucleus is short compared with the electron's
orbital period (\emph{A.~B. Migdal}, 1941; \emph{J. Levinger}, 1953).

\SolutionHeading After the $\alpha$ particle has emerged, the perturbation
acting upon the electron is adiabatic. The effect sought is therefore
determined chiefly by times close to the nonadiabatic ``switching instant,''
when the $\alpha$ particle, already outside the nucleus and moving freely, is
still at distances small compared with the radius of the $K$ orbit. The
ionizing perturbation $V$ is then the departure of the combined nuclear and
$\alpha$-particle field from the pure Coulomb field $Z/r$. For particles of
atomic weights 4 and $A-4$, charges 2 and $Z-2$, and separation $vt$ (where
$v$ is the relative speed of the nucleus and $\alpha$ particle), the dipole
moment is
\[
  \frac{2(A-4)-(Z-2)4}{A}vt=\frac{2(A-2Z)}A\,vt.
\]
Therefore, the dipole term of the field of the nucleus and the $\alpha$-particle
is given by\footnote{If the difference $A-2Z$ is small, it may also be necessary
to include the next, quadrupole term.}

\[
  V=\frac{2(A-2Z)}A\,vt\frac z{r^3},
\]
where the $z$ axis points along $\mathbf v$. The matrix element of this
perturbation can be reduced to that of $z$. Taking the matrix element of the
electron's equation of motion $\ddot z=-Zz/r^3$, we obtain
\[
  \left(\frac z{r^3}\right)_{0k}=\frac{(E-E_0)^2}{Z}z_{0k}.
\]
\SourcePage{190}{193}
By (41.2), the transition probability sought for either of the two $K$-shell
electrons is
\[
  dw=2\left|\int_0^\infty V_{0k}e^{i(E_0-E)t}\,dt\right|^2dk
  =\frac{8(A-2Z)^2v^2}{A^2Z^2}|z_{0k}|^2\frac{dk}{2\pi}.
\]
To evaluate the integral, an additional damping factor $e^{-\lambda t}$,
$\lambda>0$, is inserted in the integrand, and $\lambda\to0$ is then taken in
the result. For the matrix element of $z=r\cos\theta$, note that the initial
orbital angular momentum is $l=0$, so $\cos\theta$ has a nonzero matrix
element only for transition to a state with $l=1$. In that case
\[
  |(\cos\theta)_{01}|^2=\frac13,\qquad
  |z_{0k}|^2=\frac13|r_{0k}|^2.
\]
Evaluation of $r_{0k}$ using the radial functions $R_{00}$ and $R_{k1}$ then
gives
\[
  dw=\frac{2^{11}(A-2Z)^2v^2}{3A^2Z^6(1+k^2/Z^2)^5}
  f\left(\frac kZ\right)k\,dk,
\]
where $f$ is the function defined in Problem 4.
